\section{Beta-function area-row derivation}
\label{app:area}

\S\ref{sec:constraints} states the inscribed-area constraint row as a
closed-form Beta function without deriving it; this appendix supplies the
algebra (source: \dossierref{\S3.4}).
The per-surface area under Eq.~\eqref{eq:surfaces} is
\begin{equation}
  \mathcal{A}_{\text{surf}} = \int_0^1 \zeta(\psi)\,d\psi
  = \int_0^1 \Big[C(\psi)\sum_i A_i S_i(\psi) + \psi\,\zeta_T\Big]\,d\psi
  = \sum_i A_i \int_0^1 C(\psi)S_i(\psi)\,d\psi + \frac{\zeta_T}{2},
\end{equation}
using linearity of the integral (itself a consequence of
Eq.~\eqref{eq:surfaces}'s linearity in $A$, \S\ref{sec:cst-linearity}) and
$\int_0^1\psi\,d\psi=1/2$ for the trailing-edge term. The remaining integral,
per basis function $S_i(\psi)=K_i\psi^i(1-\psi)^{n-i}$ with
$K_i=\binom{n}{i}$ and $C(\psi)=\psi^{N_1}(1-\psi)^{N_2}$, is
\begin{align}
  \int_0^1 C(\psi)S_i(\psi)\,d\psi
    &= K_i \int_0^1 \psi^{N_1}(1-\psi)^{N_2}\,\psi^i(1-\psi)^{n-i}\,d\psi \\
    &= K_i \int_0^1 \psi^{\,i+N_1}(1-\psi)^{\,n-i+N_2}\,d\psi.
\end{align}
The Euler integral of the first kind is, by definition,
$B(p,q) = \int_0^1 t^{p-1}(1-t)^{q-1}\,dt$ for $\Re(p),\Re(q)>0$. Matching
exponents, $p-1 = i+N_1 \Rightarrow p = i+N_1+1$ and
$q-1 = n-i+N_2 \Rightarrow q = n-i+N_2+1$, giving exactly
Eq.~\eqref{eq:area-beta}:
\begin{equation}
  \int_0^1 C(\psi)S_i(\psi)\,d\psi = K_i\,B(i+N_1+1,\;n-i+N_2+1).
\end{equation}
For the default round-nose/sharp-tail class ($N_1=0.5$, $N_2=1.0$), this is
$K_i\,B(i+1.5,\,n-i+2)$, always well-defined since $i+1.5>0$ and
$n-i+2>0$ for every $i\in\{0,\dots,n\}$, so no term in the sum is singular
despite the class function's own $\sqrt\psi$ non-analyticity at $\psi=0$
(\S\ref{sec:constraints}'s point that the LE singularity affects flow-
residual \emph{rows}, not CST's own geometric \emph{functionals}, which stay
perfectly regular). Stacking one such row per coefficient
gives the full area constraint row $g$ in $g\cdot A = b$, with $b$ the
target inscribed area (or, for the T7/T8 self-consistency protocol, the
area of the CST fit itself, per the target-consistent-RHS discipline
$b=g\cdot A^\ast$ established at T4: the constraint right-hand side must
always be built from what the reference geometry \emph{actually has}, never
from an idealised or rounded value, or the constraint row itself introduces a
spurious residual floor no Newton iteration can remove).

\section{DOF accounting worked example (T7 configuration)}
\label{app:dof}

Eq.~\eqref{eq:dof} ($M+K=n_{A,\text{free}}+n_\alpha$) is an assertion, not a
derived quantity (\S\ref{sec:newton-system}); this appendix works it for the
T7 winning configuration and states the general rule it specialises.

\paragraph{The T7 case.} $n=8$ per side $\Rightarrow$ 9 coefficients per
surface, 18 total ($A_{u,0..8}$, $A_{l,0..8}$). \texttt{le\_treatment=
prescribed} removes 2 of these from the free set: $A_{u,0}$ and $A_{l,0}$
are geometrically fixed to the reference fit's own values, not Newton
unknowns (\S\ref{sec:falsifiable}), leaving $n_{A,\text{free}}=16$. Because
$A_{u,0}/A_{l,0}$ are prescribed directly
rather than tied by an equality row, no separate shared-LE constraint row is
needed for this configuration: $K=0$. $\alpha$ is fixed for the
self-consistency protocol (removing the camber--$\alpha$ null direction,
\S\ref{sec:ablations}): $n_\alpha=0$. Eq.~\eqref{eq:dof} therefore requires
$M = 16 + 0 - 0 = 16$ target-station equations, exactly the QR-pivoted
station count reported throughout \S\ref{sec:falsifiable}--\S\ref{sec:ablations}
(\pth{experiments/results/t7_naca2412/diagnostics.json}, \texttt{static.
dof\_accounting}: \texttt{n\_A\_free=16, M=16, K=0}).

\paragraph{The general rule.} Restating Eq.~\eqref{eq:dof} in full generality:
for $n$ Bernstein coefficients per surface (2 surfaces, so $2(n+1)$ raw
coefficients), $c_{LE}\in\{0,1\}$ constraint rows tying the leading-edge
coefficients (shared-LE row, used only when \texttt{le\_treatment=none} or
\texttt{lem}, not when \texttt{prescribed} pins the coefficients directly),
$p_{LE}\in\{0,2\}$ coefficients removed by prescription (2 when
\texttt{le\_treatment=prescribed}, else 0), and any additional geometric
constraint rows $K_{\text{extra}}$ (TE thickness, inscribed area, curvature,
\S\ref{sec:constraints}):
\begin{equation}
  n_{A,\text{free}} = 2(n+1) - p_{LE}, \qquad K = c_{LE} + K_{\text{extra}},
  \qquad M = n_{A,\text{free}} + n_\alpha - K.
\end{equation}
This is the Stage-1 (isolated-airfoil) specialisation of the dossier's
original rigid-body bookkeeping
(\S\ref{sec:newton-system}'s $\mathrm{DOF}=2(n+1)-1_{\text{shared LE}}-
1_{\text{TE fixed}}+3_{\text{scale/rotate/translate}}+1_{\text{stagger or
}\alpha}$): the three rigid overall modes are not exposed as free Newton
unknowns in this implementation (chord, position and rotation are fixed by
convention, not solved for), so they contribute neither to
$n_{A,\text{free}}$ nor to $M$, and Eq.~\eqref{eq:dof} is the resulting
simpler, directly-assertable statement. The rule's lineage is Lighthill's
own well-posedness bookkeeping~\citep{lighthill1945new}: a conformal-mapping
inverse solve is well-posed only when the number of free shape parameters
exactly matches the number of independent closure conditions (contour
closure, LE radius, circulation) the target and geometry family jointly
admit. Eq.~\eqref{eq:dof} is that same counting argument, discretised for
a finite-dimensional CST basis and asserted programmatically rather than
checked by hand.

\section{The onset-closure Jacobian rows}
\label{app:onset}

\S\ref{sec:transition} describes \texttt{\_residual\_transition\_forced}
(ADR-0003) narratively; this appendix states its Jacobian-row structure
explicitly, since it is the one residual block in this paper's extended
system that is not simply reused unchanged from mfoil
(\S\ref{sec:newton-system}). At a forced-transition onset station with
endpoint states $U_1$ (upstream, still formally laminar-indexed) and $U_2$
(downstream, freshly turbulent), the vendor three-row station residual is
$R=[R_{\text{mom}},\,R_{\text{shape}},\,R_{\text{lag}}]$
(\texttt{mfoil.py}, \texttt{residual\_station}). Rows 0--1
(momentum, shape parameter) are vendor's own ordinary turbulent closure
evaluated directly across $[x_1,x_2]$: an unchanged code path, so their
$U$- and $x$-Jacobian blocks are exactly mfoil's own
$\partial R_{\text{mom},\text{shape}}/\partial U_{1,2}$,
$\partial R_{\text{mom},\text{shape}}/\partial x$, reused verbatim. Row 2
(shear-lag) is replaced with the onset condition
$U_2[2] - c_{ttr}(U_2) = 0$, where $c_{ttr}$ is the same transition-
correlation root-shear-stress function (\texttt{get\_cttr}) vendor's own
\texttt{update\_transition} uses to initialise a newly-turbulent node.
Because $c_{ttr}$ is a function of $U_2$ (specifically $\theta,H_k,Re_\theta,
u_e$ at the downstream node) only, its Jacobian row has the structure
\begin{equation}
  R_U[2,\,0{:}4] = 0 \quad(\text{no } U_1 \text{ dependence}), \qquad
  R_U[2,\,4{:}8] = e_{sa} - \left.\frac{\partial c_{ttr}}{\partial U_2}\right|_{U_2},
  \qquad R_x[2,\,:] = 0,
\end{equation}
where $e_{sa}=[0,0,1,0]$ selects the shear-lag component of $U_2$ and
$\partial c_{ttr}/\partial U_2$ is \texttt{get\_cttr}'s own analytic
derivative block (\texttt{cttr\_U}), already computed by vendor code for the
identical purpose inside \texttt{update\_transition}: reused, not
re-derived. The row's independence from $U_1$ and from $x$ is not an
approximation: it follows structurally from $c_{ttr}$'s own functional
form, which was verified by direct inspection of \texttt{get\_cttr}
(ADR-0003, ``Modeling choice''), and it is what makes this replacement
row well-defined at a station where the natural (unforced) closure's
internal $x_t$-solve has no solution (\S\ref{sec:transition}).

\section{Reproducibility statement}
\label{app:repro}

\begin{itemize}
\item \textbf{Repository.} \url{https://github.com/SharathSPhD/CINS.git}
  (branch \texttt{main} at every gate closure, per this project's workflow
  discipline). \pth{vendor/mfoil/mfoil.py}~\citep{fidkowski2022coupled}
  is vendored unmodified (MIT licence) and is never edited; all solver
  access is through \pth{src/cins/solver/mfoil_adapter.py}.
\item \textbf{Manifest scheme.} Every experiment run in this paper writes
  an immutable manifest under \pth{experiments/results/} recording the
  git SHA, a hash of the fully-resolved \pth{configs/*.yaml} used, the
  random seed (\texttt{experiment.seed}, \texttt{42} throughout), and
  enough metadata to regenerate the run byte-for-byte
  (\citep{cinsrepo} condition~4). No number in this paper is
  hand-edited; every table cell traces to a \texttt{result.json} or
  \texttt{diagnostics.json} artifact under \pth{experiments/results/},
  cited inline by path throughout \S\ref{sec:falsifiable}--
  \S\ref{sec:results}.
\item \textbf{One-command regeneration.} T8 ablation-matrix figures (D-6
  overlay, $n$-sweep, flow-solve bar chart, station-selection comparison):
  \texttt{.venv/bin/python -m cins.benchmarks.figures}. This paper's new
  geometry/Cp/flow-solution/H1-gallery figures
  (\S\ref{sec:falsifiable}, \S\ref{sec:h1-panel}): \texttt{.venv/bin/python
  -m cins.benchmarks.paper\_figures}. Both re-run the underlying
  monolithic solves fresh rather than reading cached scalars, since neither
  \texttt{t7\_naca2412/diagnostics.json} nor any \texttt{t8/*/result.json}
  stores the CST coefficient vectors or field distributions themselves
  (only scalar summaries; verified directly against the on-disk schema
  before writing \texttt{paper\_figures.py}, not assumed). Individual T7/T8
  cells: \texttt{.venv/bin/python experiments/run\_t7.py} and
  \texttt{.venv/bin/python -m cins.benchmarks.runner} respectively. Tests:
  \texttt{python -m pytest tests/ -q}.
\item \textbf{Gate-closure evidence.} \citep{cinsrepo},
  \texttt{T3-closure.md}, \texttt{T4-T7-closure.md} record the six-condition
  gate-closure contract (\citep{cinsrepo}) for each corresponding
  task, including adversarial-review findings and their resolution;
  \citep{cinsrepo} through \texttt{ADR-0004} record the four
  binding architectural decisions this paper's Formulation section depends
  on (scipy shim, vendor NACA5 shim, forced-transition shim, realisability
  semantics).
\end{itemize}

\bibliographystyle{plainnat}
\bibliography{p1_refs}
